% SGA 3, Expose V, Loop 1 English reconstruction.
% Scope: Section 6 opening and quasi-section reduction through the groupoid V_*.
% Authority: Polo--Gille Exp5-13oct24.pdf, component-local p. 20
% (combined-reader p. 282), re-edited page header 270.
% Jacob Reinhold's CC BY 4.0 Markdown is comparison/drafting material only.

\section*{6. 存在拟截面时的取商}
\phantomsection\label{sec:sga3-V-6}
\addcontentsline{toc}{section}{6. 存在拟截面时的取商}

现在证明一个技术性引理, 将用于我们所要证明的两个定理.
设 \(S\) 是概形, 并设
\SGAThreeDiagram[0.54\textwidth]
  {figures/exp5/p270_section6_groupoid_X.png}
  {群胚 \(X_\ast\).}
  {fig:sga3-V-6-X-groupoid}
  {第 V 篇, 分篇第~20 页 / 合订本第~282 页 / 重版本页眉页码~270}
是一个 \((\mathbf{Sch}/S)\)-群胚.
群胚 \(X_\ast\) 的\emph{拟截面 (quasi-section)} 是指
\(X_0\) 的任意满足下列两个条件的子概形 \(U\):
\begin{enumerate}
\item[(1)]\phantomsection\label{def:sga3-V-6-quasisection-1}%
  \(d_1\) 在 \(d_0^{-1}(U)\) 上的限制 \(v\)
  是从 \(d_0^{-1}(U)\) 到 \(X_0\) 的有限、局部自由、满射态射.
\item[(2)]\phantomsection\label{def:sga3-V-6-quasisection-2}%
  \(U\) 的每个由关于 \(X_\ast\) 所定义的等价关系
  (\hyperref[par:sga3-V-3-e]{\S3.e)}) 两两等价的点组成的子集 \(E\),
  都包含在 \(U\) 的一个仿射开集中.
  \footnote{\label{note:sga3-V-31}编者注: 若 \(x,y\in E\),
  则有 \(z\in X_1\), 使 \(d_1(z)=x\) 且 \(d_0(z)=y\);
  因而 \(z\) 属于 \(v^{-1}(x)\), 后者由
  \hyperref[def:sga3-V-6-quasisection-1]{(1)} 为有限集.
  所以 \(E\) 包含在有限集
  \(d_0(v^{-1}(x))=d_0d_1^{-1}(x)\cap U\) 中.
  源 PDF 在第一个等式中印作 \(d_1(x)=x\);
  这里采用类型正确的读法 \(d_1(z)=x\).}
\end{enumerate}

若 \(U\) 是 \(X_\ast\) 的拟截面, 则由 \(X_\ast\)
和包含映射 \(U\hookrightarrow X_0\) 诱导
(见 \hyperref[par:sga3-V-3-a]{\S3.a)}) 的
\((\mathbf{Sch}/S)\)-群胚
\SGAThreeDiagram[0.48\textwidth]
  {figures/exp5/p270_section6_groupoid_U.png}
  {由 \(X_\ast\) 诱导的群胚 \(U_\ast\).}
  {fig:sga3-V-6-U-groupoid}
  {第 V 篇, 分篇第~20 页 / 合订本第~282 页 / 重版本页眉页码~270}
满足定理~\hyperref[thm:sga3-V-4.1]{4.1} 的假设. 事实上, 置
\[
  V=d_0^{-1}(U),
\]
并用 \(u\) 和 \(v\) 分别表示由 \(d_0\) 和 \(d_1\)
诱导的以 \(V\) 为源的态射:
\SGAThreeDiagram[0.43\textwidth]
  {figures/exp5/p270_section6_span_X0_V_U.png}
  {态射 \(v:V\to X_0\) 和 \(u:V\to U\).}
  {fig:sga3-V-6-span-X0-V-U}
  {第 V 篇, 分篇第~20 页 / 合订本第~282 页 / 重版本页眉页码~270}

由 \hyperref[par:sga3-V-3-b]{\S3.b)}, 有笛卡儿方块
\SGAThreeDiagram[0.27\textwidth]
  {figures/exp5/p270_section6_cartesian_square.png}
  {定义 \(U_1\) 的笛卡儿方块.}
  {fig:sga3-V-6-U1-cartesian}
  {第 V 篇, 分篇第~20 页 / 合订本第~282 页 / 重版本页眉页码~270}
故由\hyperref[def:sga3-V-6-quasisection-1]{条件~(1)},
\(u_1\) 满射且有限局部自由. 因此,
\hyperref[def:sga3-V-6-quasisection-1]{条件~(1)} 连同
\hyperref[def:sga3-V-6-quasisection-2]{条件~(2)},
保证群胚 \(U_\ast\) 满足定理~\hyperref[thm:sga3-V-4.1]{4.1} 的假设.
特别地, \(\operatorname{Coker}(u_0,u_1)\) 在 \((\mathbf{Sch}/S)\) 中存在.
此外, \(d_0\) 有截面, 故 \(u\) 是普遍有效满态射
(见第 IV 篇, 1.12).
\footnote{法文重版本误引为第 III 篇, 该篇并无 \S1.12;
见 Canon 记录 \texttt{SGA-SRC-000856}.}
命题~\hyperref[prop:sga3-V-3.1]{3.1} 于是表明,
\(\operatorname{Coker}(u_0,u_1)\) 与群胚 \(V_\ast\)
的余核 \(\operatorname{Coker}(v_0,v_1)\) 相同:
\SGAThreeDiagram[0.48\textwidth]
  {figures/exp5/p270_section6_groupoid_V.png}
  {群胚 \(V_\ast\).}
  {fig:sga3-V-6-V-groupoid}
  {第 V 篇, 分篇第~20 页 / 合订本第~282 页 / 重版本页眉页码~270}

% Translation coverage: entire assigned component; adjoining components retain the original statement/proof seams.
% Source lines: 82; translated body lines before this audit comment: 83.
% Math preserved: 50 inline formulas, 1 display blocks, 0 equation environments; mathematical prose translated where present.
% Labels: 4/4; hyperrefs: 10/10; SGARefLinks: 0/0; inline targets: 0/0; footnotes: 2/2.
% Figure arguments preserved; all rendered captions and natural-language source locators translated.
% Unresolved TODOs: 0; compilation and visual inspection remain with the volume integrator.
